\documentclass[11pt]{article}
\usepackage[a4paper,margin=2.6cm]{geometry}
\usepackage{amsmath,amssymb,amsthm}
\newtheorem{lemma}{Lemma}
\newtheorem{proposition}[lemma]{Proposition}
\theoremstyle{remark}\newtheorem*{remark}{Remark}
\title{Two analytic statements that do not move RH:\\
\large visibility at $\gamma_1$, and why the $2$-adic mass sits at $2^{\pm1}$}
\author{D.\ Joubert}
\date{3 September 2026}
\begin{document}
\maketitle
\begin{abstract}
\noindent
Two calculations. (1)~An off-line zero at height $\gamma_1$ makes the
windowed form negative for every $\sigma\neq0$ small, because the
App.~B term $\hat f(\rho)\hat f(1-\rho)$ opens at order $-\sigma^2$
on the ground state while $\lambda_{\min}$ is exponentially small.
(2)~The difference of Fourier transforms $F_S-F_\infty$ is a sum of
dilations by $2^{\mathbb Z}$; conjugating the scaling $\vartheta(\lambda)$
against those dilations puts the extra trace mass at $\lambda=2^{\pm1}$,
with leading weight the local Weil factor $(\log2)/\sqrt2$. Neither
statement is RH.
\end{abstract}

\section{Visibility at $\gamma_1$}
Let $\tau=L/2$, $f$ even on $[-\tau,\tau]$, and
\[
\hat f(\tfrac12+\sigma+i\gamma)=\int_{-\tau}^{\tau}f(x)\,e^{\sigma x}e^{i\gamma x}\,dx.
\]
App.~B of \emph{quorum-theorem}: $\hat g(\rho)=\hat f(\rho)\,\hat f(1-\rho)$
for $g=f\star\tilde f$. Expand at fixed $\gamma$ in $\sigma$,
\[
\hat f(\tfrac12+\sigma+i\gamma)=F_0+\sigma F_1+\tfrac12\sigma^2 F_2+O(\sigma^3),
\]
with $F_0=\hat f(\tfrac12+i\gamma)$ real (even $f$) and $F_1=\int f(x)\,x\,e^{i\gamma x}\,dx=iM$, $M=\int f(x)\,x\sin(\gamma x)\,dx$ real: for even $f$ the first moment is purely imaginary.
Then
\[
\hat g(\rho)=F_0^2-\sigma^2 |F_1|^2+O(\sigma^2 F_0)+O(\sigma^4)=F_0^2-\sigma^2M^2+\ldots,
\]
Take $f$ the ground state of $Q_L$ on $V_N$. Hyper-nullity
(\emph{zeros-from-the-radical}) gives $|F_0(\gamma_1)|\approx C_\gamma\lambda_{\min}$,
exponentially small. The moment $|F_1(\gamma_1)|$ is $O(\tau\|f\|)$ a priori; it depends on $N$: on $V_9$ at $\mu=11$ it is $O(1)$ ($M^2\approx0.54$ from $\hat g=-5.4\times10^{-5}$ at $\sigma=0.01$), while on $V_{47}$ at the same $\mu$ the same quantity, $|\partial_\gamma\hat v_0(\gamma_1)|$, is $3.5\times10^{-3}$ and decays like $e^{-0.6\gamma}$ along the band (notebook \S76): the visibility scale in $\sigma$ moves by two orders of magnitude between the two bases. Hence on the $V_9$ vector
\[
Q_L(f)+\hat g(\rho)=\lambda_{\min}-\sigma^2 |F_1(\gamma_1)|^2+O(\sigma^2\lambda_{\min})+O(\sigma^4).
\]
\begin{lemma}[visibility at $\gamma_1$, small $\sigma$]
If $F_1(\gamma_1)\ne0$, then $Q_L+\hat g(\rho)$ is negative for every
sufficiently small $\sigma\neq0$.
\end{lemma}
The hypothesis $F_1(\gamma_1)\ne0$ is a non-vanishing of one Fourier
moment of an explicit eigenvector. It is measured, not proved for every
$N,\mu$. It is the only analytic gap in the lemma.

\paragraph{High zeros.}
$|\hat\eta_0(\tfrac12+\sigma+i\gamma)|\le e^{|\sigma|\tau}/(|\gamma|\sqrt L)$
and similarly for $n\ge1$ away from $\omega_n$. The extra Gram at height
$\gamma$ has operator norm $O(e^{2|\sigma|\tau}/\gamma^2)$ on $V_N$.
At $\gamma=80$, $\mu=11$, this is $10^{-3}$ or less and does not move
$\lambda_{\min}$ (measured: stays $10^{-22}$). Visibility is local in
height: a window of type $\tau$ sees an off-line zero near $\gamma_1$,
not one at height $\gg 2\pi N/L$. A single certified $Q_{11}\succ0$ is
compatible with an off-line zero at $10^6$. That is Weil on a compact
window, not RH.

\section{Why the mass sits at $2^{\pm1}$}
On the slice $\mathrm{ord}_2=0$,
\[
F_S=\tfrac12\Bigl(\sum_{n\ge0}M_{2^n}-M_{1/2}\Bigr)F_\infty,
\]
where $(M_a\hat g)(\rho)=\hat g(a\rho)$. Thus
\[
F_S-F_\infty=\tfrac12\sum_{n\ge1}M_{2^n}F_\infty-\tfrac12 M_{1/2}F_\infty-\tfrac12 F_\infty.
\]
The scaling $\vartheta(\lambda)g(r)=\lambda^{-1/2}g(r/\lambda)$
conjugates dilations of the frequency:
$M_a\,\vartheta(\lambda)\,M_a^{-1}=\vartheta(\lambda/a)$ up to the
normalization of $d^*\lambda$. Every extra copy $M_{2^n}$ in the trace
$\mathrm{Tr}(\hat P_\Lambda P_\Lambda\vartheta(\lambda))$ produces a
translate of the archimedean peak by $\lambda\mapsto 2^{\pm n}\lambda$.
The first shells $n=\pm1$ sit at $\lambda=2$ and $\lambda=\tfrac12$.
Higher shells $2^{\pm n}$ are present in the formula and not resolved
as isolated peaks at $\Lambda\le8$.

The coefficient of the first shell is the $k=1$ term of the local
Weil factor at $p=2$:
$(\log2)\,2^{-1/2}=(\log2)/\sqrt2$.
At $\Lambda=4$ and $\Lambda=8$ the hard-window mass of
$\tau_S-\tau_\infty$ around $\lambda=2$ extrapolates to that number
as the cell size $h\to0$ ($0.488$ and $0.491$). That is a numerical
limit of an explicit distribution, not a new identity. The
archimedean peak at $\lambda=1$ and the factor $2h(1)\log\Lambda$
are the other two terms of Connes (1999), Theorem~4; they are not
re-derived here.

\begin{remark}
The hard window around $\lambda=\tfrac12$ at $\Lambda=8$ overshoots
($w=1.17$ at $96$ cells per unit): it swallows mass that is not the
first shell. The peak at $\lambda=2$ is clean. An analytic shell
expansion would separate $n=\pm1$ from $n=\pm2,\pm3$ without a window.
That calculation is not done.
\end{remark}

\section*{Status}
Lemma~1 is a Taylor argument plus one measured non-vanishing.
Proposition on the locus $2^{\pm1}$ is conjugation of dilations;
the weight is identified with the local factor and confirmed by
extrapolation, not proved as an operator coefficient.
No movement of RH. The two remaining analytic targets of the
repository --- $\delta_p\le e^{-7p}$ and $\kappa_p\ge e^{-4p}$ on
the bottom vector, uniformly in $\mu$ --- are not touched.
\end{document}
